\begin{figure}[t]
\centering
\resizebox{\columnwidth}{!}{%
\begin{tikzpicture}
\begin{axis}[
  width=8.4cm, height=5.7cm,
  ybar, bar width=12pt,
  symbolic x coords={GCN, Heat, QI},
  xtick=data,
  ymin=0, ymax=0.24,
  ylabel={error (lower is better)},
  legend pos=north east, legend cell align=left,
  ymajorgrids, grid style={gray!18},
  tick label style={font=\footnotesize}, label style={font=\small},
  legend style={font=\footnotesize, fill=white, fill opacity=0.85, draw=gray!40},
  axis line style={gray!50},
  error bars/y dir=both, error bars/y explicit,
  enlarge x limits=0.28,
]
\addplot[draw=steel!80!black, fill=steel!35] coordinates {
  (GCN,0.163) +- (0,0.016)
  (Heat,0.069) +- (0,0.005)
  (QI,0.071) +- (0,0.009)};
\addlegendentry{MAE}
\addplot[draw=amber!80!black, fill=palegold] coordinates {
  (GCN,0.202) (Heat,0.085) (QI,0.077)};
\addlegendentry{far-node AURC}
\end{axis}
\end{tikzpicture}}
\caption{Exp~B, operator comparison on the LEO node-potential task ($264$ satellites,
$5$ seeds, param-matched). Global operators (Heat, QI) beat the local GCN by roughly
$\globalgain\times$. Heat and the quantum-inspired (QI) walk tie on MAE within the seed
spread (error bars); QI shows a small, theory-consistent edge only on the far-node AURC.}
\label{fig:expB}
\end{figure}
